\chapter{信号分析}
\label{chap:signal-analysis}

第~\ref{chap:graph-theory}章刚把局部关系写成图，并用图拉普拉斯谱描述节点信号的
变化模式。更早的线性代数把矩阵理解为线性映射，矩阵分析又说明了怎样用谱坐标观察算子。
当输入带有时间、空间位置等顺序时，还会出现一类额外结构：同一个局部作用可以沿位置
重复使用。直接把这种作用写成一般矩阵当然可行，但那会隐藏重复的系数，也无法立即看出
边界、频率和采样怎样改变结果。

本章以两个有限序列
\begin{equation}
  \bm{x}=(1,2,1),\qquad
  \bm{h}=(1,-1)
  \label{eq:signal-analysis-running-example}
\end{equation}
作为贯穿例子。我们先从平移不变的线性滤波推出卷积，再把序列系数写成多项式；
复指数提供另一组坐标，使循环卷积化为逐点乘法。快速傅里叶变换、Toom--Cook和
Winograd类算法都利用“先改变表示、再做较少乘法、最后恢复结果”的思路，但它们
使用的分解并不相同。采样部分说明什么时候这种表示仍保留原信号，图信号和状态核则
把规则序列上的结论分别推广到不规则关系和递推系统。

\section{离散信号与线性滤波}
\label{sec:signal-analysis-discrete-filtering}

\subsection{信号是带索引的数值对象}

\term{离散时间信号}（discrete-time signal）是从整数索引到实数或复数的函数，
\[
  x:\mathbb{Z}\to\mathbb{C},\qquad t\mapsto x[t].
\]
索引$t$可以表示时间，也可以表示一维空间位置。信号的含义来自索引顺序，而不只是
一组数值：交换$x[0]$与$x[1]$通常会改变局部变化和滤波输出。若只有有限多个位置
非零，就称信号具有有限支撑。式~\eqref{eq:signal-analysis-running-example}中的
$\bm{x}$表示$x[0]=1,x[1]=2,x[2]=1$，其余位置取零；$\bm{h}$也作同样的零延拓。

对整数$r$，定义平移算子
\[
  (T_r x)[t]=x[t-r].
\]
$r>0$时，原来位于$t$的值移动到$t+r$。这个符号约定很重要，因为另一种常见约定
$x[t+r]$会把平移方向反过来。单位脉冲
\[
  \delta[t]=
  \begin{cases}
    1,&t=0,\\
    0,&t\neq0
  \end{cases}
\]
满足$T_r\delta[t]=\delta[t-r]$。任意有限支撑信号都可分解为平移脉冲之和：
\begin{equation}
  x[t]=\sum_{r\in\mathbb{Z}}x[r]\delta[t-r].
  \label{eq:signal-analysis-impulse-expansion}
\end{equation}
对每个固定$t$，只有$r=t$的一项非零。这个简单分解将把一般输入的响应归结为系统对
单个脉冲的响应。

\subsection{线性与平移不变共同推出卷积}

一个离散系统$\mathcal{H}$把输入信号$x$映射为输出信号$y=\mathcal{H}x$。
若对任意标量$a,b$及输入$x,z$都有
\[
  \mathcal{H}(ax+bz)=a\mathcal{H}x+b\mathcal{H}z,
\]
就称系统为线性的。若
\[
  \mathcal{H}(T_r x)=T_r(\mathcal{H}x)
  \qquad\text{对每个 }r\in\mathbb{Z}\text{成立},
\]
就称系统具有\term{平移不变性}（shift invariance）。这表示输入整体移动后，
输出只作同样移动，系统在不同位置使用同一规则。

令$h=\mathcal{H}\delta$，称$h$为系统的\term{脉冲响应}（impulse response）。
若$x$具有有限支撑，式~\eqref{eq:signal-analysis-impulse-expansion}是有限和，
线性性足以把$\mathcal{H}$移入求和。对无穷序列，记
\[
  \ell_1(\mathbb{Z})
  =
  \left\{
    x:\mathbb{Z}\to\mathbb{C}
    \,\middle|\,
    \sum_{t\in\mathbb{Z}}|x[t]|<\infty
  \right\},
  \qquad
  \lVert x\rVert_1=\sum_{t\in\mathbb{Z}}|x[t]|.
\]
也就是说，$\ell_1(\mathbb{Z})$由绝对可和序列组成。对一般的
$x\in\ell_1(\mathbb{Z})$，还需假设
$\mathcal{H}:\ell_1(\mathbb{Z})\to\ell_1(\mathbb{Z})$是有界线性且平移不变的
算子；这里有界表示存在与$x$无关的常数$C$，使
\[
  \lVert\mathcal{H}x\rVert_1\leq C\lVert x\rVert_1.
\]
令
\[
  x^{(N)}=\sum_{|r|\leq N}x[r]T_r\delta,
\]
则$x^{(N)}\to x$在$\ell_1$范数下成立；有界性蕴含连续性，因而
$\mathcal{H}x^{(N)}\to\mathcal{H}x$。同时
$h=\mathcal{H}\delta\in\ell_1(\mathbb{Z})$，而且
\[
  \sum_{t\in\mathbb{Z}}\sum_{r\in\mathbb{Z}}
  |x[r]h[t-r]|
  =
  \lVert x\rVert_1\lVert h\rVert_1
  <\infty,
\]
所以下面的无穷和绝对收敛。于是，
有限支撑情形直接由线性性得到、$\ell_1$情形通过上述极限得到
\begin{align}
  (\mathcal{H}x)[t]
  &=
  \mathcal{H}\left(\sum_{r\in\mathbb{Z}}x[r]T_r\delta\right)[t]
  \notag\\
  &=
  \sum_{r\in\mathbb{Z}}x[r](T_rh)[t]
  =
  \sum_{r\in\mathbb{Z}}x[r]h[t-r].
  \label{eq:signal-analysis-lti-convolution}
\end{align}

式~\eqref{eq:signal-analysis-lti-convolution}定义\term{离散卷积}
（discrete convolution），记为
\begin{equation}
  (x*h)[t]=\sum_{r\in\mathbb{Z}}x[r]h[t-r].
  \label{eq:signal-analysis-linear-convolution}
\end{equation}
核索引$t-r$随$r$增加而反向读取，这个翻转是数学卷积定义的一部分。卷积满足交换律；
令$s=t-r$便有
\[
  \sum_r x[r]h[t-r]
  =
  \sum_s h[s]x[t-s]
  =
  (h*x)[t].
\]
在绝对可和等允许交换求和次序的条件下，它还满足结合律和分配律。

\subsection{相关与卷积只差一个方向，却回答不同问题}

\term{互相关}（cross-correlation）常写为
\begin{equation}
  (x\star h)[t]
  =
  \sum_{r\in\mathbb{Z}}x[t+r]\overline{h[r]}.
  \label{eq:signal-analysis-cross-correlation}
\end{equation}
它把$h$按原方向滑过$x$，并在复数情形取共轭；上划线表示保持实部不变、把虚部
符号反转。若定义
$\widetilde{h}[r]=\overline{h[-r]}$，则
\[
  (x\star h)[t]=(x*\widetilde{h})[t].
\]
因此相关可以用卷积实现，但二者的核含义不同。信号处理中，相关常用于比较一个模板与
不同位置的局部片段；许多神经网络库把不翻转核的相关运算称为“卷积”。只要训练和
推理始终采用同一约定，参数可以适应该方向；在手算、复用滤波器或推导伴随算子时，
仍必须明确是否翻转以及是否取共轭。

\subsection{线性卷积与循环卷积采用不同边界}

设$x$的支撑为$0,\ldots,m-1$，$h$的支撑为$0,\ldots,r-1$。零延拓下的线性卷积
只可能在$0,\ldots,m+r-2$非零，因此完整输出长度为$m+r-1$。对贯穿例，
\begin{align*}
  (x*h)[0]&=1,\\
  (x*h)[1]&=2-1=1,\\
  (x*h)[2]&=1-2=-1,\\
  (x*h)[3]&=-1.
\end{align*}
所以
\begin{equation}
  \bm{x}*\bm{h}=(1,1,-1,-1).
  \label{eq:signal-analysis-running-linear-convolution}
\end{equation}
核$(1,-1)$在这里计算当前值减去前一个值；开头和末尾的结果还包含零延拓带来的边界
跳变。

若信号被视为长度$N$的周期序列，索引应按模$N$回绕。长度$N$的
\term{循环卷积}（circular convolution）定义为
\begin{equation}
  (x*_N h)[t]
  =
  \sum_{r=0}^{N-1}x[r]h[(t-r)\bmod N],
  \qquad t=0,\ldots,N-1.
  \label{eq:signal-analysis-circular-convolution}
\end{equation}
把例中的$\bm{h}$补成$(1,-1,0)$并取$N=3$，可得
\[
  \bm{x}*_3\bm{h}=(0,1,-1).
\]
它等于把式~\eqref{eq:signal-analysis-running-linear-convolution}中索引相差$3$
的项相加：首项$1$与末项$-1$叠在一起。循环边界不是零边界的近似，而是另一个线性
算子。

固定有限长度后，线性卷积可写成Toeplitz矩阵乘法，循环卷积可写成循环矩阵乘法。
例如
\[
  \begin{bmatrix}
    1&0&0\\
    -1&1&0\\
    0&-1&1\\
    0&0&-1
  \end{bmatrix}
  \begin{bmatrix}1\\2\\1\end{bmatrix}
  =
  \begin{bmatrix}1\\1\\-1\\-1\end{bmatrix}.
\]
矩阵的重复对角线来自参数共享。裁剪为“valid”输出、补零后保持输入长度，或使用反射、
复制等边界，都会改变矩阵的首尾行。平移不变性在无限网格或循环群上最为精确；
有限区间上的边界处理通常会破坏全局平移不变性。

“valid”并不是另一种卷积定义，而是只保留核完全落在已知输入区间内的输出。按本章
$h[0],\ldots,h[r-1]$的因果索引约定，当$m\geq r$时，这些输出对应
$t=r-1,\ldots,m-1$，共$m-r+1$项；若$m<r$，则没有这样的输出。补零得到的首尾
输出仍有定义，但它们使用了人为指定的区间外数值。因而比较两个实现时，除了核系数，
还要同时核对延拓方式、保留的输出索引和核的方向。

\section{多项式乘法、求值与余数}
\label{sec:signal-analysis-polynomial-arithmetic}

\subsection{卷积系数就是多项式乘积系数}

有限序列可以用形式多项式编码。对
\[
  \bm{x}=(x_0,\ldots,x_{m-1}),\qquad
  \bm{h}=(h_0,\ldots,h_{r-1}),
\]
定义
\[
  X(z)=\sum_{i=0}^{m-1}x_i z^i,\qquad
  H(z)=\sum_{j=0}^{r-1}h_j z^j.
\]
变量$z$在这里首先是记录索引的形式符号。乘积为
\begin{align}
  X(z)H(z)
  &=
  \sum_{i=0}^{m-1}\sum_{j=0}^{r-1}x_i h_jz^{i+j}
  \notag\\
  &=
  \sum_{t=0}^{m+r-2}
  \left(\sum_{i+j=t}x_i h_j\right)z^t.
  \label{eq:signal-analysis-polynomial-convolution}
\end{align}
括号中的系数正是$(x*h)[t]$。卷积的交换律和结合律也可分别从多项式乘法的交换律和
结合律得到。

贯穿例对应
\[
  X(z)=1+2z+z^2,\qquad H(z)=1-z,
\]
因而
\begin{equation}
  P(z)=X(z)H(z)=1+z-z^2-z^3.
  \label{eq:signal-analysis-running-polynomial-product}
\end{equation}
其四个系数与式~\eqref{eq:signal-analysis-running-linear-convolution}完全一致。
这种对应没有减少计算，只是把“哪些下标相加”改写成“哪些幂次相乘”。真正的计算收益
来自下一步：多项式除了系数表示，还可以用函数值或余数表示。

\subsection{求值与插值是同一多项式的另一组坐标}

次数小于$L$的多项式空间维数为$L$。选取$L$个两两不同的数
$a_0,\ldots,a_{L-1}$，求值映射
\[
  E(P)=\bigl(P(a_0),\ldots,P(a_{L-1})\bigr)
\]
是一个线性映射。若$E(P)=\bm{0}$，则$P$有$L$个不同根；次数小于$L$的非零
多项式至多有$L-1$个根，所以只能有$P=0$。求值映射是同维有限空间之间的单射，
因而可逆。其逆就是\term{多项式插值}（polynomial interpolation）。

具体地，Lagrange基多项式
\[
  \ell_j(z)=
  \prod_{\substack{0\leq k<L\\k\neq j}}
  \frac{z-a_k}{a_j-a_k}
\]
满足$\ell_j(a_k)=\delta_{j,k}$。因此
\begin{equation}
  P(z)=\sum_{j=0}^{L-1}P(a_j)\ell_j(z),
  \qquad \deg P<L.
  \label{eq:signal-analysis-lagrange-interpolation}
\end{equation}
若要计算$P=XH$，可以先分别求$X(a_j)$与$H(a_j)$，再用
$P(a_j)=X(a_j)H(a_j)$逐点相乘，最后插值。逐点乘法数是$L$，代价被转移到求值
与插值中的线性组合。选点、问题规模和数值精度决定这种交换是否值得。

\subsection{带余除法把高次信息压到指定模中}

求值是余数的一种特殊情形。为建立更一般的表示，先考虑多项式除法。

\begin{theorem}[多项式带余除法]
\label{thm:signal-analysis-polynomial-division}
设$F$为一个域，$A(z),M(z)\in F[z]$且$M\neq0$。存在唯一的多项式
$Q(z),R(z)$使
\[
  A(z)=Q(z)M(z)+R(z),
  \qquad \deg R<\deg M.
\]
\end{theorem}

\begin{proof}
若$\deg A<\deg M$，取$Q=0,R=A$即可。否则设$A$与$M$的最高次项分别为
$a_pz^p$与$m_qz^q$，其中$p\geq q$。因为$F$是域，$m_q$可逆；从$A$中减去
$(a_p/m_q)z^{p-q}M$便消去最高次项。重复此过程，次数严格下降，有限步后得到次数
小于$\deg M$的余数。

若另有$A=Q'M+R'$且$\deg R',\deg R<\deg M$，则
$(Q-Q')M=R'-R$。左侧若非零，次数至少为$\deg M$；右侧若非零，次数小于
$\deg M$，矛盾。因此$Q=Q'$且$R=R'$。
\end{proof}

带余除法还给出后文重构所需的互素判据。对非零多项式$A,B$反复作
\[
  A=Q_0B+R_1,\qquad
  B=Q_1R_1+R_2,\qquad \ldots
\]
每个非零余数的次数都严格下降，所以过程有限步终止；最后一个非零余数经过首项归一化，
就是$A,B$的\term{最大公因式}（greatest common divisor）。从最后一个等式向前
回代，每个余数都能写成$A,B$的多项式线性组合。因此存在$U,V\in F[z]$使
\begin{equation}
  U(z)A(z)+V(z)B(z)=\gcd(A,B).
  \label{eq:signal-analysis-polynomial-bezout}
\end{equation}
这就是多项式的Bézout恒等式。特别地，$A,B$互素当且仅当可找到$U,V$使左侧等于
$1$；“互素”在这里表示最大公因式为非零常数。

例如
\[
  z^2+1=(z+1)(z-1)+2,
\]
所以$z-1$与$z^2+1$互素，而且回代立即给出
\[
  1=\frac{1}{2}(z^2+1)-\frac{1}{2}(z+1)(z-1).
\]
后文构造中国剩余重构系数时，使用的正是这一步，而不是假设逆元凭空存在。

余数记为$A\bmod M$，并把相差$M$倍数的多项式视为同一个
\term{剩余类}（residue class）。当$M(z)=z-a$时，带余除法给出
\[
  A(z)=Q(z)(z-a)+A(a),
\]
所以关于一次多项式$z-a$的余数就是点值$A(a)$。

模多项式还编码边界。因为在模$z^N-1$的剩余类中有$z^N=1$，
$z^{N+s}$与$z^s$的系数会合并。于是
\begin{equation}
  \left(X(z)H(z)\right)\bmod(z^N-1)
  =
  \sum_{t=0}^{N-1}(x*_Nh)[t]z^t.
  \label{eq:signal-analysis-polynomial-circular-convolution}
\end{equation}
线性卷积与循环卷积的差别由此变成“保留普通乘积”还是“再对$z^N-1$取模”。
若$\deg(XH)<N$，取模不会合并任何系数，两种卷积在补零后的前$N$项相同；
若次数达到$N$，高次项就会回绕。

\section{复指数与频率表示}
\label{sec:signal-analysis-frequency-representation}

\subsection{复指数把平移变成相位乘法}

后文会同时使用复数、复向量和复矩阵。记$\mathrm{i}^2=-1$；对
$z=a+\mathrm{i}b$，其\term{复共轭}（complex conjugate）与模长分别为
\[
  \overline z=a-\mathrm{i}b,
  \qquad
  |z|=\sqrt{z\overline z}.
\]
当$z\neq0$时，可写成$z=|z|e^{\mathrm{i}\arg z}$，其中$\arg z$是相角。
互相关式~\eqref{eq:signal-analysis-cross-correlation}中的上划线也表示复共轭。
对$\bm{u},\bm{v}\in\mathbb{C}^N$，本章采用复内积
\[
  \langle\bm{u},\bm{v}\rangle
  =
  \sum_{t=0}^{N-1}u[t]\overline{v[t]}.
\]
在第二个变量上取共轭使
$\langle\bm{u},\bm{u}\rangle=\sum_t|u[t]|^2$为非负实数。复矩阵
$\bm{A}$的\term{共轭转置}（conjugate transpose）记为
$\bm{A}^*=\overline{\bm{A}}^{\top}$；若$\bm{A}^*\bm{A}=\bm{I}$，则称
$\bm{A}$为酉矩阵。

时域中的平移会改动许多坐标。复指数
\[
  \phi_{\omega}[t]=e^{\mathrm{i}\omega t}
\]
却在平移后只差一个标量：
\[
  (T_r\phi_{\omega})[t]
  =
  e^{\mathrm{i}\omega(t-r)}
  =
  e^{-\mathrm{i}\omega r}\phi_{\omega}[t].
\]
因此复指数是平移算子的特征函数，频率$\omega$标记不同的特征模式，系数
$e^{-\mathrm{i}\omega r}$记录平移产生的相位变化。这解释了频率表示为何适合分析
平移不变系统：所有平移共享同一组模式。

\subsection{离散傅里叶变换建立有限维正交坐标}

对长度$N$的序列，令
\[
  \omega_N=e^{-2\pi\mathrm{i}/N}.
\]
\term{离散傅里叶变换}（discrete Fourier transform, DFT）及逆变换定义为
\begin{align}
  \widehat{x}[k]
  &=
  \sum_{t=0}^{N-1}x[t]\omega_N^{kt},
  &&k=0,\ldots,N-1,
  \label{eq:signal-analysis-dft}\\
  x[t]
  &=
  \frac{1}{N}\sum_{k=0}^{N-1}
  \widehat{x}[k]\omega_N^{-kt},
  &&t=0,\ldots,N-1.
  \label{eq:signal-analysis-idft}
\end{align}
正变换也就是在$N$次单位根$1,\omega_N,\ldots,\omega_N^{N-1}$处计算多项式
$X(z)$的值。它既是频率坐标，也是第~\ref{sec:signal-analysis-polynomial-arithmetic}
节求值表示的一种特殊选点。

逆变换来自单位根的正交关系。对整数$t,s$，
\begin{equation}
  \sum_{k=0}^{N-1}\omega_N^{k(t-s)}
  =
  \begin{cases}
    N,&t\equiv s\pmod N,\\
    0,&t\not\equiv s\pmod N.
  \end{cases}
  \label{eq:signal-analysis-root-orthogonality}
\end{equation}
第一种情形每项都是$1$。第二种情形令$q=\omega_N^{t-s}\neq1$，利用$q^N=1$，
几何级数之和为$(1-q^N)/(1-q)=0$。把式~\eqref{eq:signal-analysis-dft}
代入式~\eqref{eq:signal-analysis-idft}，再交换两个有限求和，只有$s=t$的项留下，
从而恢复$x[t]$。这也证明DFT没有丢失长度$N$序列的信息。

若把DFT矩阵记为$\bm{F}$，其中$F_{k,t}=\omega_N^{kt}$，则
\[
  \bm{F}^{*}\bm{F}=N\bm{I}.
\]
因此$\bm{F}/\sqrt{N}$是酉矩阵。第~\ref{chap:matrix-analysis}章的谱语言在这里
给出一个直接解释：$\bm{F}/\sqrt{N}$才是保持二范数的正交坐标变换。本章定义的
未归一化DFT满足
$\lVert\widehat{\bm{x}}\rVert_2=\sqrt{N}\lVert\bm{x}\rVert_2$，
逆变换中的$1/N$会补回这一尺度。

\subsection{幅度、相位与实信号的共轭对称}

每个非零复系数都可写成
\[
  \widehat{x}[k]=
  |\widehat{x}[k]|e^{\mathrm{i}\arg\widehat{x}[k]}.
\]
模长描述该频率坐标的大小，相角描述它相对于所选时间原点的相位。若把信号循环平移$r$
个位置，DFT变为
\[
  \widehat{T_rx}[k]
  =
  \omega_N^{kr}\widehat{x}[k].
\]
幅度不变，相位随$k$线性改变。仅保留幅度会丢失位置信息；两个具有相同幅度谱的信号
不必具有相同波形。

若$x[t]$为实数，则
\begin{align*}
  \overline{\widehat{x}[k]}
  &=
  \sum_{t=0}^{N-1}x[t]\omega_N^{-kt}\\
  &=
  \sum_{t=0}^{N-1}x[t]\omega_N^{(N-k)t}
  =
  \widehat{x}[(N-k)\bmod N].
\end{align*}
这称为\term{共轭对称}（conjugate symmetry）。频率$0$总是实数；$N$为偶数时，
$k=N/2$对应的Nyquist频率也为实数。其余频率成共轭对出现，所以实信号的完整复频谱
包含冗余，但直接丢弃一半频率时仍要保留正确的缩放和端点约定。

\subsection{Parseval关系说明能量只是在坐标间重排}

由$\bm{F}^{*}\bm{F}=N\bm{I}$，
\[
  \lVert\widehat{\bm{x}}\rVert_2^2
  =
  \bm{x}^{*}\bm{F}^{*}\bm{F}\bm{x}
  =
  N\lVert\bm{x}\rVert_2^2.
\]
于是得到离散的\term{Parseval关系}
\begin{equation}
  \sum_{t=0}^{N-1}|x[t]|^2
  =
  \frac{1}{N}\sum_{k=0}^{N-1}
  |\widehat{x}[k]|^2.
  \label{eq:signal-analysis-parseval}
\end{equation}
这里的$1/N$来自本章把归一化全部放在逆变换中。若正逆变换各用$1/\sqrt{N}$，
等式两侧便没有额外因子。比较不同资料中的频谱数值前，必须先核对变换的符号与归一化
约定。

\subsection{连续周期信号需要区分收敛方式}

有限维DFT只涉及有限求和，不存在无穷级数的收敛问题。把同一正交展开思路推广到
连续周期信号后，才需要区分平方平均收敛与逐点收敛。对周期为$T$的平方可积连续信号
$f(t)$，复指数$e^{\mathrm{i}2\pi kt/T}$在内积
\[
  \langle f,g\rangle
  =
  \frac{1}{T}\int_0^T f(t)\overline{g(t)}\,\mathrm{d}t
\]
下正交。其\term{傅里叶系数}（Fourier coefficient）为
\[
  c_k=
  \frac{1}{T}\int_0^T
  f(t)e^{-\mathrm{i}2\pi kt/T}\,\mathrm{d}t.
\]
形式上，傅里叶级数写为
\begin{equation}
  f(t)\sim\sum_{k\in\mathbb{Z}}c_k
  e^{\mathrm{i}2\pi kt/T}.
  \label{eq:signal-analysis-fourier-series}
\end{equation}
若只假设$f\in L^2[0,T]$，等号应理解为平方平均意义下收敛；逐点收敛还需要额外
规则性，并且跳跃点处通常收敛到左右极限的平均。频率展开不是无条件的逐点恒等式。
在相同归一化下，连续周期信号还满足Parseval关系
\begin{equation}
  \frac{1}{T}\int_0^T|f(t)|^2\,\mathrm{d}t
  =
  \sum_{k\in\mathbb{Z}}|c_k|^2.
  \label{eq:signal-analysis-continuous-parseval}
\end{equation}
它是$L^2$中的能量恒等式，不要求傅里叶级数在每一点都等于$f$。

一个具体边界例子是周期为$2\pi$的方波：在$0<t<\pi$取$1$，在
$-\pi<t<0$取$-1$，再作周期延拓。直接积分得到$c_0=0$，偶数$k$对应
$c_k=0$，奇数$k$对应
\[
  c_k=\frac{2}{\pi\mathrm{i}k}.
\]
方波在除跳跃点外的每个连续点由级数恢复，在跳跃点则收敛到左右极限的平均$0$；
与此同时，式~\eqref{eq:signal-analysis-continuous-parseval}仍在$L^2$意义下成立。
这一区别说明，平均平方误差趋于零并不能逐点排除跳跃附近的振荡。

\section{卷积定理与快速变换}
\label{sec:signal-analysis-fast-transforms}

\subsection{循环卷积在频域中逐点相乘}

\begin{theorem}[离散循环卷积定理]
\label{thm:signal-analysis-circular-convolution-theorem}
设$x,h$为长度$N$的复序列，DFT采用
式~\eqref{eq:signal-analysis-dft}的约定，则
\[
  \widehat{x*_Nh}[k]
  =
  \widehat{x}[k]\widehat{h}[k],
  \qquad k=0,\ldots,N-1.
\]
\end{theorem}

\begin{proof}
从定义出发，并把$a=(t-r)\bmod N$作为新的求和索引：
\begin{align*}
  \widehat{x*_Nh}[k]
  &=
  \sum_{t=0}^{N-1}
  \sum_{r=0}^{N-1}
  x[r]h[(t-r)\bmod N]\omega_N^{kt}\\
  &=
  \sum_{r=0}^{N-1}
  \sum_{a=0}^{N-1}
  x[r]h[a]\omega_N^{k(r+a)}\\
  &=
  \left(\sum_{r=0}^{N-1}x[r]\omega_N^{kr}\right)
  \left(\sum_{a=0}^{N-1}h[a]\omega_N^{ka}\right).
\end{align*}
第二个等号中，$t$遍历一个完整剩余类时$a$也恰好遍历$0,\ldots,N-1$；
即使$t$与$r+a$相差$N$的整数倍，$\omega_N^{kN}=1$也保证幂值不变。最后两项
分别是$\widehat{x}[k]$和$\widehat{h}[k]$。
\end{proof}

因此
\begin{equation}
  x*_Nh
  =
  \mathcal{F}_N^{-1}
  \left(
    \mathcal{F}_N(x)\odot\mathcal{F}_N(h)
  \right).
  \label{eq:signal-analysis-frequency-domain-convolution}
\end{equation}
循环卷积矩阵由DFT基对角化，每个频率上的“特征值”就是核的DFT系数。这个结论同时
连接了第~\ref{chap:matrix-analysis}章的谱分解和本章的多项式求值：循环矩阵、
模$z^N-1$的乘法与单位根求值
描述的是同一个算子。

若要得到长度$m+r-1$的线性卷积，必须把两个序列都补零到
\begin{equation}
  N\geq m+r-1.
  \label{eq:signal-analysis-linear-padding-condition}
\end{equation}
此时乘积次数小于$N$，模$z^N-1$不会发生系数回绕。只把两者补到
$N\geq\max(m,r)$不够；它保证输入装得下，却未保证乘积装得下。

\subsection{FFT利用单位根的递归结构}

直接计算$N$个DFT系数，每个需要$N$项求和，运算量为$O(N^2)$。
\term{快速傅里叶变换}（fast Fourier transform, FFT）不是另一种变换，而是一族
利用求值点结构快速计算DFT的算法\scite{cooley1965fft}以$N$为偶数的基二
Cooley--Tukey分解为例，
将输入的偶数项与奇数项分别组成
\[
  x_{\mathrm{e}}[r]=x[2r],\qquad
  x_{\mathrm{o}}[r]=x[2r+1].
\]
记它们的长度$N/2$ DFT为$E[k]$和$O[k]$。因为$\omega_N^2=\omega_{N/2}$，
对$0\leq k<N/2$有
\begin{align}
  \widehat{x}[k]
  &=
  E[k]+\omega_N^kO[k],
  \notag\\
  \widehat{x}[k+N/2]
  &=
  E[k]-\omega_N^kO[k].
  \label{eq:signal-analysis-fft-butterfly}
\end{align}
第二式利用了$\omega_N^{N/2}=-1$。两个较小DFT的同一对结果通过一次加法和一次
减法产生两个频率，并没有舍弃后一半频率。

当$N=2^q$时，递推
\[
  T(N)=2T(N/2)+O(N)
\]
给出$T(N)=O(N\log N)$。其他合数长度可按不同因子分解，质数长度也有相应算法；
“使用FFT”不等于必须补到二的幂。补到更方便的长度可能降低实现常数，也可能增加
无用频点和存储，需按具体规模判断。

对贯穿例补零到$N=4$：
\[
  \bm{x}_4=(1,2,1,0),\qquad
  \bm{h}_4=(1,-1,0,0).
\]
对$\bm{x}_4$，偶数分支$(1,1)$的二点DFT为$(2,0)$，奇数分支$(2,0)$的
二点DFT为$(2,2)$。取$\omega_4=-\mathrm{i}$并按
式~\eqref{eq:signal-analysis-fft-butterfly}合并，得到
\[
  (2+2,\ 0-2\mathrm{i},\ 2-2,\ 0+2\mathrm{i})
  =
  (4,-2\mathrm{i},0,2\mathrm{i}).
\]
同理，$\bm{h}_4$的偶、奇分支DFT分别为$(1,1)$与$(-1,-1)$，蝶形合并得到
$(0,1+\mathrm{i},2,1-\mathrm{i})$。这一步显示四点结果如何由两个二点变换复用，
而不只是直接代入DFT定义。

在$1,-\mathrm{i},-1,\mathrm{i}$四个单位根处，
\[
  \widehat{\bm{x}}_4=(4,-2\mathrm{i},0,2\mathrm{i}),\qquad
  \widehat{\bm{h}}_4=(0,1+\mathrm{i},2,1-\mathrm{i}).
\]
例如逆变换的四个分子依次为
\[
  4-2\mathrm{i}+2\mathrm{i},\quad
  4+2+2,\quad
  4+2\mathrm{i}-2\mathrm{i},\quad
  4-2-2,
\]
除以$4$后恢复$(1,2,1,0)$。这一步同时核对了频率顺序、指数符号和$1/N$归一化，
任何一项约定改变都要相应修改逆变换。

逐点乘积为
\[
  (0,2-2\mathrm{i},0,2+2\mathrm{i}),
\]
逆DFT恢复$(1,1,-1,-1)$。若只用$N=3$，末项会回绕到首项，得到的便是
$(0,1,-1)$而非线性卷积。

\subsection{Toom--Cook用一般求值点减少系数乘法}

Toom--Cook方法把多项式分块，选择足够多的求值点，逐点相乘后再插值。对本例的
二次多项式$X$与一次多项式$H$，乘积次数至多为$3$，可使用
$0,1,-1,\infty$四个点；其中“在$\infty$处求值”只表示读取最高次项系数，
并非把无穷大传给数值程序\scite{math-knuth1997seminumerical}记
\[
  r_0=X(0)H(0),\quad
  r_+=X(1)H(1),\quad
  r_-=X(-1)H(-1),\quad
  r_\infty=x_2h_1.
\]
若$P(z)=c_0+c_1z+c_2z^2+c_3z^3$，则
\begin{align}
  c_0&=r_0,\qquad c_3=r_\infty,\notag\\
  c_1&=\frac{r_+-r_-}{2}-c_3,\notag\\
  c_2&=\frac{r_++r_-}{2}-c_0.
  \label{eq:signal-analysis-toom-reconstruction}
\end{align}
这些式子由$P(1)$与$P(-1)$的和、差直接解出。

对式~\eqref{eq:signal-analysis-running-polynomial-product}，
\[
  (r_0,r_+,r_-,r_\infty)=(1,0,0,-1),
\]
代入式~\eqref{eq:signal-analysis-toom-reconstruction}得到
$(c_0,c_1,c_2,c_3)=(1,1,-1,-1)$。直接展开需要$3\times2=6$个输入系数与
核系数之间的乘法，这个求值方案只需四个点值乘法，但新增了求值、插值和常数乘除。
对长多项式递归分块时才得到相应的渐近改进；对短核，只数核心乘法会掩盖变换和访存
成本。

\subsection{余数分解与多项式中国剩余定理}

点值只保留关于一次多项式的余数。Winograd类算法更一般地选择若干互素模多项式，
分别计算低次余数，再重构所需乘积或滤波输出。其代数依据是下面的多项式中国剩余
定理\scite{math-winograd1980arithmetic}

\begin{theorem}[多项式中国剩余定理]
\label{thm:signal-analysis-polynomial-crt}
设$F$为一个域，$m_1(z),\ldots,m_s(z)\in F[z]$两两互素，并令
$M(z)=\prod_{j=1}^{s}m_j(z)$。映射
\[
  [A]_M
  \longmapsto
  \bigl([A]_{m_1},\ldots,[A]_{m_s}\bigr)
\]
是环同构。换言之，关于各$m_j$的余数唯一确定次数小于$\deg M$的
$A\bmod M$。
\end{theorem}

\begin{proof}
若两个多项式$A,B$在每个$m_j$下余数相同，则每个$m_j$都整除$A-B$。
这里需要说明两两互素如何推出乘积整除。若互素多项式$a,b$都整除$q$，写成
$q=ac$；由Bézout恒等式$ua+vb=1$可得
\[
  c=uac+vbc=uq+vbc,
\]
故$b$整除$c$，进而$ab$整除$q$。两两互素时，任意若干因子的乘积仍与其余因子
互素，反复应用这个Euclid引理便得$M$整除$A-B$。所以$A,B$属于同一个模$M$
剩余类，映射为单射。

为构造任意给定余数$r_j$的原像，令$M_j=M/m_j$。由于$M_j$与$m_j$互素，
Bézout恒等式给出多项式$u_j,v_j$使
\[
  u_jM_j+v_jm_j=1.
\]
于是$e_j=u_jM_j$关于$m_j$的余数为$1$，关于其余$m_k$的余数为$0$。
多项式
\[
  A=\sum_{j=1}^{s}r_je_j
\]
便具有全部指定余数。再对$M$取余即可得到唯一的次数小于$\deg M$代表元。
\end{proof}

仍用同一个乘积，并取
\[
  M(z)=z^4-1=(z-1)(z+1)(z^2+1).
\]
这三个因子在$\mathbb{R}[z]$中两两互素，总次数为$4$。分别计算$XH$的余数：
\[
  P(1)=0,\qquad P(-1)=0,
\]
而在模$z^2+1$下有$z^2=-1$，
\[
  X(z)\equiv2z,\qquad
  H(z)\equiv1-z,\qquad
  X(z)H(z)\equiv2+2z.
\]
设重构结果$R$的次数小于$4$。前两个余数为零说明$R$可写为
\[
  R(z)=(z^2-1)(az+b).
\]
再对$z^2+1$取模，得到$R\equiv-2az-2b$。与$2+2z$比较可得
$a=b=-1$，所以
\[
  R(z)=1+z-z^2-z^3.
\]
由于原乘积次数已经小于$4$，这个模$z^4-1$的代表元就是完整乘积本身。
从输入输出形状看，这次余数计算把长度$3$与长度$2$的系数向量映射为长度$4$的完整
线性卷积。若只要求中间若干项，Winograd最小滤波会据此重新选择余数与恢复映射；
不能把这里恢复全部四项的变换直接当作任意输出块的最小算法。

\begin{example}[可复核的Winograd \(F(2,3)\)输出块]
  \label{ex:signal-analysis-winograd-f23}
  考虑一维“valid”互相关。长度为\(4\)的输入块
  \(\bm d=(d_0,d_1,d_2,d_3)^{\top}\in\mathbb{R}^{4}\)与长度为\(3\)的核
  \(\bm g=(g_0,g_1,g_2)^{\top}\in\mathbb{R}^{3}\)产生长度为\(2\)的输出块
  \[
    y_0=d_0g_0+d_1g_1+d_2g_2,
    \qquad
    y_1=d_1g_0+d_2g_1+d_3g_2.
  \]
  \(F(2,3)\)选取关于\(z\)、\(z-1\)、\(z+1\)的余数，并补上最高次系数坐标；
  前三项分别对应在\(0,1,-1\)处的值。核侧变换直接是这些余数坐标的缩放，输入侧
  还要针对只恢复上述两个中间系数作固定线性组合。按本例的互相关方向，一组变换为
  \begin{equation}
    \bm y
    =
    \bm A^{\top}
    \left[
      (\bm G\bm g)\odot(\bm B^{\top}\bm d)
    \right],
    \label{eq:signal-analysis-winograd-f23}
  \end{equation}
  其中
  \begin{align*}
    \bm B^{\top}
    &=
    \begin{bmatrix}
      1&0&-1&0\\
      0&1&1&0\\
      0&-1&1&0\\
      0&1&0&-1
    \end{bmatrix},\\
    \bm G
    &=
    \begin{bmatrix}
      1&0&0\\
      1/2&1/2&1/2\\
      1/2&-1/2&1/2\\
      0&0&1
    \end{bmatrix},\\
    \bm A^{\top}
    &=
    \begin{bmatrix}
      1&1&1&0\\
      0&1&-1&-1
    \end{bmatrix}.
  \end{align*}
  因而形状依次为
  \[
    \mathbb{R}^{4}\xrightarrow{\bm B^{\top}}\mathbb{R}^{4},
    \qquad
    \mathbb{R}^{3}\xrightarrow{\bm G}\mathbb{R}^{4},
    \qquad
    \mathbb{R}^{4}\xrightarrow{\bm A^{\top}}\mathbb{R}^{2}.
  \]
  取\(\bm d=(1,2,3,4)^{\top}\)、\(\bm g=(2,-1,1)^{\top}\)，可逐步核对
  \[
    \bm B^{\top}\bm d=(-2,5,1,-2)^{\top},
    \qquad
    \bm G\bm g=(2,1,2,1)^{\top}.
  \]
  四次逐坐标乘法得到
  \(\bm m=(-4,5,2,-2)^{\top}\)，再作输出变换便有
  \[
    \bm A^{\top}\bm m=(3,5)^{\top}.
  \]
  直接互相关也给出
  \(y_0=2-2+3=3\)与\(y_1=4-3+4=5\)。这个小块把直接计算的六次
  输入--核乘法换成四次逐坐标乘法；矩阵变换中的加减和常数乘法仍须计入实现成本。
\end{example}

该例也揭示FFT、Toom--Cook与余数算法的联系和区别。FFT在$z^N-1$的全部复单位根
处求值，并利用根的规则结构快速变换；Toom--Cook可选一般求值点并插值完整乘积；
Winograd类构造从更一般的互素余数出发，常以减少特定小块滤波所需的乘法为目标。
一次因子的CRT与普通点值插值重合，但高次模因子会保留一个低次多项式余数，不能再把
它说成单个点值。

\subsection{算法比较必须同时核对算术、误差与边界}

长度相近的两个$n$项序列直接卷积需要$O(n^2)$次乘加；FFT路线在适合的变换长度上
需要$O(N\log N)$次变换运算和$N$次频域乘法。Toom--Cook通过递归分块减少组合
乘法，Winograd最小滤波则针对固定输出块减少核心乘法。渐近量级或乘法次数都不是完整的
运行时间模型：加减、常数乘法、复数运算、数据变换、缓存复用和目标硬件同样重要。

数值误差也依赖表示。归一化DFT矩阵是酉矩阵，坐标变换本身条件良好，但有限精度的
多层蝶形仍会累积舍入误差；Toom--Cook与Winograd的求值、插值矩阵可能含尺度差异
较大的系数，相近量相减也可能放大相对误差。整数或有限域中，若模数和变换长度允许，
可以用数论变换避免浮点舍入，但还要防止模运算后的系数范围歧义。不存在脱离数据类型、
规模和误差目标而普遍最优的卷积算法。

边界条件更不能由快速算法自动决定。FFT默认循环卷积，线性卷积必须满足
式~\eqref{eq:signal-analysis-linear-padding-condition}；Toom--Cook恢复的是所设
多项式乘积，网络互相关还需翻转或重新排列核；Winograd小块算法只恢复设计时指定的
输出，步长、空洞率或边界变化后需要重新推导。不同算法可以计算同一映射，但前提是
输入延拓、核方向、输出裁剪与数值精度均保持一致。

\section{采样与信息损失}
\label{sec:signal-analysis-sampling}

\subsection{一个混叠例子说明丢失为何不可逆}

离散序列常来自连续信号$f(t)$在间隔$T_s$上的观测：
\[
  x[n]=f(nT_s),\qquad \nu_s=\frac{1}{T_s}.
\]
这里$\nu_s$是采样频率。若不限制原信号允许包含的频率，采样通常不能唯一确定
连续信号。令$T_s=1$，考虑
\[
  f_1(t)=\cos(0.8\pi t),\qquad
  f_2(t)=\cos(1.2\pi t).
\]
两者具有不同频率，但在整数样本上
\[
  f_2[n]
  =
  \cos(1.2\pi n)
  =
  \cos(2\pi n-0.8\pi n)
  =
  \cos(0.8\pi n)
  =
  f_1[n].
\]
只看样本，无法判断原信号来自$0.8\pi$还是$1.2\pi$。任何后续插值器都只能在多个
相容连续信号中作选择，不能从相同数据恢复已经丢失的区别。

这个例子也解释了为什么“先下采样，再学习恢复”没有一般的信息保证。模型可以利用
训练分布偏好猜测常见的高频细节，但若两个输入在采样后完全相同，任何确定性后续映射
对它们也只能给出同一输出。恢复效果来自额外先验或数据规律，而不是下采样本身可逆。
上面的纯余弦在整条实轴上不属于$L^2(\mathbb{R})$，这里只用来直观展示功率信号中的
采样歧义；下面转入平方可积带限信号，给出函数空间闭合的重构结论。

\subsection{采样会把连续频谱周期复制}

设$f\in L^2(\mathbb{R})$。当$f\in L^1(\mathbb{R})\cap L^2(\mathbb{R})$时，
连续时间傅里叶变换先由普通积分定义为
\[
  F(\Omega)=\int_{-\infty}^{\infty}
  f(t)e^{-\mathrm{i}\Omega t}\,\mathrm{d}t.
\]
Plancherel定理把这个变换唯一延拓到$L^2(\mathbb{R})$；本节的$F$均指这一
$L^2$意义下的傅里叶变换。若
$F(\Omega)=0$对几乎处处的$|\Omega|>\Omega_{\max}$成立，就称$f$带限于
$\Omega_{\max}$。\term{Paley--Wiener带限空间}就是满足该条件的
$L^2(\mathbb{R})$函数所组成的空间。

由Plancherel定理，$F\in L^2(\mathbb{R})$。带限条件又把其支撑限制在有限测度
区间内，所以Cauchy--Schwarz不等式给出
\[
  \lVert F\rVert_1
  \leq
  \sqrt{2\Omega_{\max}}\,\lVert F\rVert_2
  <\infty.
\]
因此逆变换积分
\[
  f_{\mathrm{c}}(t)
  =
  \frac{1}{2\pi}\int_{-\Omega_{\max}}^{\Omega_{\max}}
  F(\Omega)e^{\mathrm{i}\Omega t}\,\mathrm{d}\Omega
\]
定义了$f$的一个连续代表元。下面仍把这个代表元记为$f$，点值$f(nT_s)$与采样都
针对它定义。

令采样角频率为
\[
  \Omega_s=\frac{2\pi}{T_s},
\]
并把区间
$I_s=[-\pi/T_s,\pi/T_s]$称为一个\term{基带}（baseband）。把频谱副本按间隔
$\Omega_s$相加，定义
\[
  P_F(\Omega)
  =
  \frac{1}{T_s}
  \sum_{k\in\mathbb{Z}}F(\Omega-k\Omega_s).
\]
由于$F$带限，对几乎每个固定$\Omega$只有有限项非零；在$I_s$上只会出现有限多个
$F$的平移，因此$P_F\in L^2(I_s)$。函数$P_F$以$\Omega_s$为周期，并且它在
$I_s$上的第$n$个傅里叶系数为
\begin{align*}
  \frac{1}{\Omega_s}
  \int_{I_s}
  P_F(\Omega)e^{\mathrm{i}\Omega nT_s}\,\mathrm{d}\Omega
  &=
  \frac{1}{2\pi}
  \int_{\mathbb{R}}
  F(\Omega)e^{\mathrm{i}\Omega nT_s}\,\mathrm{d}\Omega\\
  &=
  f(nT_s)
  =
  x[n].
\end{align*}
第一步把各个平移区间拼回整条实轴，并使用
$e^{\mathrm{i}nT_s\Omega_s}=e^{\mathrm{i}2\pi n}=1$；第二步是逆傅里叶变换。
由$P_F\in L^2(I_s)$及傅里叶级数的Parseval关系，其系数序列
$\{x[n]\}_{n\in\mathbb{Z}}$属于$\ell_2(\mathbb{Z})$。因此离散样本的频率表示
可以定义为$L^2(I_s)$意义下的傅里叶级数
\[
  X_s(\Omega)
  =
  \sum_{n\in\mathbb{Z}}
  x[n]e^{-\mathrm{i}\Omega nT_s}.
\]
它与$P_F$具有相同的傅里叶系数，故二者在$L^2(I_s)$中相同，即
\begin{equation}
  X_s(\Omega)
  =
  \frac{1}{T_s}
  \sum_{k\in\mathbb{Z}}F(\Omega-k\Omega_s).
  \label{eq:signal-analysis-spectral-replication}
\end{equation}
采样不是只截取一份频谱，而是把原频谱按$\Omega_s$周期复制并相加。

若$\Omega_s>2\Omega_{\max}$，相邻副本互不重叠，可以用理想低通滤波器取回中央
副本。若$\Omega_s<2\Omega_{\max}$，允许频带的相邻副本在正测度区间上重叠，
采样映射在整个Paley--Wiener空间上不再是单射。不同连续频率由此落到同一离散频率，
这种现象称为\term{混叠}（aliasing）。临界等号
$\Omega_s=2\Omega_{\max}$则不同：在当前$L^2$模型中，相邻副本至多在端点相接，
零测端点不改变频谱的$L^2$等价类，不能把它无条件判为混叠。

另一种常见推导把采样写成信号与Dirac脉冲列的乘积，再对脉冲列作傅里叶变换；其中的
乘积、变换和等式都要在广义函数意义下解释。本章不调用这套理论，而以上述基带傅里叶
级数作为证明接口。

\begin{theorem}[带限信号的采样重构]
\label{thm:signal-analysis-sampling}
若$f$属于上述Paley--Wiener空间，频谱支撑包含于
$[-\Omega_{\max},\Omega_{\max}]$，且$\Omega_s>2\Omega_{\max}$，
则$f$由样本$f(nT_s)$唯一确定，并可写成
\begin{equation}
  f(t)
  =
  \sum_{n\in\mathbb{Z}}
  f(nT_s)
  \operatorname{sinc}\left(\frac{t-nT_s}{T_s}\right),
  \qquad
  \operatorname{sinc}(u)=\frac{\sin(\pi u)}{\pi u}.
  \label{eq:signal-analysis-sinc-reconstruction}
\end{equation}
级数在$L^2(\mathbb{R})$中收敛，并在额外的常见正则条件下逐点或一致收敛。
\end{theorem}

\begin{proof}
记$I_s=[-\pi/T_s,\pi/T_s]$。条件$\Omega_s>2\Omega_{\max}$保证$F$的支撑
严格包含在$I_s$内，式~\eqref{eq:signal-analysis-spectral-replication}在该区间上
只有中央副本，即$F=T_sX_s$。因此$F$在$I_s$上的傅里叶系数为
$T_sf(nT_s)$，相应部分和
\[
  F_N(\Omega)
  =
  T_s\sum_{n=-N}^{N}
  f(nT_s)e^{-\mathrm{i}\Omega nT_s}
\]
在$L^2(I_s)$中收敛到$F$。把$F_N$在$I_s$外取为零，再作逆傅里叶变换。有限求和
可以直接移到积分外，得到
\begin{align*}
  f_N(t)
  &=
  \sum_{n=-N}^{N}f(nT_s)
  \frac{T_s}{2\pi}
  \int_{-\pi/T_s}^{\pi/T_s}
  e^{\mathrm{i}\Omega(t-nT_s)}\,\mathrm{d}\Omega\\
  &=
  \sum_{n=-N}^{N}f(nT_s)
  \operatorname{sinc}\left(\frac{t-nT_s}{T_s}\right).
\end{align*}
第二个等号只需计算对称区间上的指数积分；当$t=nT_s$时，按连续延拓把
$\operatorname{sinc}(0)$取为$1$。傅里叶变换在$L^2$中的等距性质说明
$F_N\to F$蕴含$f_N\to f$，于是令$N\to\infty$便得到
式~\eqref{eq:signal-analysis-sinc-reconstruction}。

若两条满足条件的信号具有相同样本，则它们频谱在$I_s$上的全部傅里叶系数相同。
傅里叶级数在$L^2(I_s)$中的唯一性说明两频谱几乎处处相同，逆变换后两信号也在
$L^2(\mathbb{R})$中相同。这同时给出唯一性。
\end{proof}

定理采用严格不等号，以便让频谱支撑严格落在基带内部并留出低通余量。在当前
Paley--Wiener $L^2$模型中，同一证明可按几乎处处意义延伸到
$\Omega_s=2\Omega_{\max}$，因为端点相接不影响$L^2$等价类。若把信号类扩大到
允许谱线的广义函数或功率信号，临界点才需要另行处理：例如
$\sin(\Omega_{\max}t)$在$T_s=\pi/\Omega_{\max}$时的全部样本均为零，但它不属于
$L^2(\mathbb{R})$。两种结论针对的信号类不同。

实际观测只有有限时长，理想sinc具有无限支撑，噪声与非带限成分也不会因定理而消失。
因此该定理给出理想信号类中的无混叠重构条件，不是有限数据上的无误差保证。

\subsection{离散时间傅里叶变换与频率响应}

采样得到的是双向无穷序列，不能直接把长度$N$的DFT当作它的频率表示。对
$x\in\ell_1(\mathbb{Z})$，定义\term{离散时间傅里叶变换}
（discrete-time Fourier transform, DTFT）
\begin{equation}
  X(e^{\mathrm{i}\omega})
  =
  \sum_{n\in\mathbb{Z}}x[n]e^{-\mathrm{i}\omega n},
  \qquad \omega\in\mathbb{R}.
  \label{eq:signal-analysis-dtft}
\end{equation}
绝对可和保证级数收敛，并使$X(e^{\mathrm{i}\omega})$成为连续的$2\pi$周期函数。
对$x\in\ell_2(\mathbb{Z})$，同一表达按傅里叶级数在
$L^2([-\pi,\pi])$中解释，等式只需几乎处处成立。DFT把$N$个数映到$N$个离散
频率坐标，并带有长度$N$的周期边界；DTFT则把无穷序列映到连续频率变量上的周期函数。

设滤波核$h$具有有限支撑或属于$\ell_1(\mathbb{Z})$。虽然复指数
$x[t]=e^{\mathrm{i}\omega t}$本身不属于$\ell_1$，$h$的绝对可和性仍保证下面的
卷积和逐点绝对收敛：
\begin{align*}
  (x*h)[t]
  &=
  \sum_r h[r]e^{\mathrm{i}\omega(t-r)}\\
  &=
  e^{\mathrm{i}\omega t}
  \sum_r h[r]e^{-\mathrm{i}\omega r}.
\end{align*}
输出仍是同一频率，只乘以
\begin{equation}
  H(e^{\mathrm{i}\omega})
  =
  \sum_r h[r]e^{-\mathrm{i}\omega r},
  \label{eq:signal-analysis-filter-frequency-response}
\end{equation}
即滤波器的\term{频率响应}（frequency response），也就是$h$的DTFT。其模长决定
该模式被放大或削弱多少，相位决定时间位置怎样变化。

例中$h=(1,-1)$的响应为
\[
  H(e^{\mathrm{i}\omega})=1-e^{-\mathrm{i}\omega},
  \qquad
  |H(e^{\mathrm{i}\omega})|=2|\sin(\omega/2)|.
\]
常量模式$\omega=0$被完全消去，接近零频的缓慢变化也被削弱，高频变化则较强。
因此该核可作一阶差分，但边界延拓仍会决定有限序列首尾出现什么响应。

\subsection{下采样前的滤波决定哪些频率被保留}

对离散信号每隔$M$个位置保留一项，
\[
  y[n]=x[Mn],
\]
称为$M$倍\term{下采样}（downsampling）。沿用
式~\eqref{eq:signal-analysis-dtft}的DTFT约定，
\begin{equation}
  Y(e^{\mathrm{i}\omega})
  =
  \frac{1}{M}
  \sum_{\ell=0}^{M-1}
  X\left(e^{\mathrm{i}(\omega+2\pi\ell)/M}\right).
  \label{eq:signal-analysis-downsampling-spectrum}
\end{equation}
要验证这个式子，把右侧的$X$按定义展开并交换有限和与级数：
\begin{align*}
  \frac{1}{M}\sum_{\ell=0}^{M-1}
  X\left(e^{\mathrm{i}(\omega+2\pi\ell)/M}\right)
  &=
  \sum_{r\in\mathbb{Z}}x[r]e^{-\mathrm{i}r\omega/M}
  \left(
    \frac{1}{M}\sum_{\ell=0}^{M-1}
    e^{-2\pi\mathrm{i}\ell r/M}
  \right)\\
  &=
  \sum_{n\in\mathbb{Z}}x[Mn]e^{-\mathrm{i}\omega n}.
\end{align*}
括号中的单位根之和只在$r$是$M$的倍数时等于$1$，其余情形为$0$，所以最后一行
正是$y[n]=x[Mn]$的频谱。对于绝对可和序列，上述交换直接成立；更一般的能量有限
序列需在$L^2$意义下解释。

输入模式$e^{\mathrm{i}\omega_0n}$经下采样后变成
$e^{\mathrm{i}M\omega_0n}$，所以输入频率映到
$M\omega_0\bmod 2\pi$。换言之，原频谱沿输出频率轴扩展$M$倍，再由$2\pi$周期
副本相加；超出主值区间的部分会回绕并可能混叠。要避免不同项重叠，应在下采样前用
低通滤波器把频率限制到$|\omega|<\pi/M$附近；这一步称为
\term{抗混叠滤波}（anti-aliasing filtering）。有限长度滤波器不能实现理想砖墙
频响，工程设计需要在过渡带、延迟、计算量和阻带衰减之间取舍。

平均池化可以看作低通滤波后下采样。长度$M$的均值核
\[
  h[n]=
  \begin{cases}
    1/M,&0\leq n<M,\\
    0,&\text{其他}
  \end{cases}
\]
具有频率响应
\begin{equation}
  H(e^{\mathrm{i}\omega})
  =
  \frac{1}{M}
  e^{-\mathrm{i}\omega(M-1)/2}
  \frac{\sin(M\omega/2)}{\sin(\omega/2)}.
  \label{eq:signal-analysis-moving-average-response}
\end{equation}
当$\omega\in2\pi\mathbb{Z}$使正弦商呈$0/0$形式时，按原有限和的连续延拓取
$H(e^{\mathrm{i}\omega})=1$。
它削弱部分高频，却不是理想低通，旁瓣仍可能在下采样后产生混叠。最大池化是非线性
运算，不能用一个固定频率响应完整描述；它同样会合并多个位置，因而通常不可逆。

可以在同一输入上比较三种处理。令
\[
  x[n]=\cos(0.2\pi n)+\cos(0.8\pi n),
\]
并作二倍下采样。直接抽取给出
\[
  x[2m]=2\cos(0.4\pi m),
\]
因为$0.8\pi$分量抽取后折叠到与$0.2\pi$分量相同的离散频率。若先用理想低通完全
删除$0.8\pi$分量，输出则为$\cos(0.4\pi m)$。若只用两点均值
$h=(1/2,1/2)$，由
\[
  \frac{\cos(\omega n)+\cos(\omega(n-1))}{2}
  =
  \cos(\omega/2)\cos\bigl(\omega n-\omega/2\bigr)
\]
可知两个分量在抽取后仍会相加，只是分别乘上$\cos(0.1\pi)$与
$\cos(0.4\pi)$并产生相位偏移。具体输出为
\[
  \cos(0.1\pi)\cos(0.4\pi m-0.1\pi)
  +
  \cos(0.4\pi)\cos(1.6\pi m-0.4\pi).
\]
第二项仍会折回较低离散频率。局部平均降低了高频项的幅度，却没有把它消除；
“先平滑”与“满足无混叠带限条件”不是同一个保证。

滤波、裁剪、量化、池化和降维都可能丢失信息，判断能否恢复必须说明限制在哪个信号类上。
若某个滤波器在一段频率上响应为零，该频段分量进入零空间；若下采样发生混叠，不同输入
落到同一输出。第~\ref{chap:linear-algebra}章的零空间和
第~\ref{chap:matrix-analysis}章的小奇异值语言因此仍适用：精确零响应意味着不可识别，
接近零的响应意味着反演对噪声敏感。

\section{图信号与状态核}
\label{sec:signal-analysis-graph-state-kernels}

\subsection{图傅里叶变换用拉普拉斯特征向量定义变化模式}

规则序列有统一的“向前一步”，平移算子因而自然；不规则关系上却没有唯一的全局平移。
设无向加权图$\mathcal{G}=(V,E)$有$n$个节点，对称非负邻接矩阵为
$\bm{W}\in\mathbb{R}^{n\times n}$，度矩阵为
\[
  \bm{D}=\operatorname{diag}(d_1,\ldots,d_n),
  \qquad d_i=\sum_jw_{i,j}.
\]
组合图拉普拉斯为$\bm{L}=\bm{D}-\bm{W}$。一个
\term{图信号}（graph signal）是给每个节点赋值的向量
$\bm{x}\in\mathbb{R}^n$。

图拉普拉斯的二次型满足
\begin{equation}
  \bm{x}^{\top}\bm{L}\bm{x}
  =
  \frac{1}{2}\sum_{i=1}^{n}\sum_{j=1}^{n}
  w_{i,j}(x_i-x_j)^2.
  \label{eq:signal-analysis-graph-dirichlet-energy}
\end{equation}
展开右侧，两个平方项分别给出$\sum_i d_ix_i^2$的两份，交叉项为
\[
  -\sum_{i=1}^{n}\sum_{j=1}^{n}w_{i,j}x_ix_j,
\]
合并后便得到左侧。因此$\bm{L}$对称半正定；
二次型越小，强连接节点上的数值越接近。

由第~\ref{chap:matrix-analysis}章的谱定理，
\[
  \bm{L}=\bm{U}\bm{\Lambda}\bm{U}^{\top},
  \qquad
  0=\lambda_0\leq\lambda_1\leq\cdots\leq\lambda_{n-1}.
\]
定义\term{图傅里叶变换}（graph Fourier transform）
\begin{equation}
  \widehat{\bm{x}}_{\mathcal{G}}
  =
  \bm{U}^{\top}\bm{x},
  \qquad
  \bm{x}
  =
  \bm{U}\widehat{\bm{x}}_{\mathcal{G}}.
  \label{eq:signal-analysis-graph-fourier}
\end{equation}
特征值小的特征向量在相邻节点间变化较缓，特征值大的模式具有较高的图Dirichlet能量，
所以可把$\lambda_k$理解为图上的频率尺度。

\subsection{规则环图把图谱接回DFT}

长度$N$的循环序列可以看作环图，每个节点与左右邻居以权重$1$相连。其拉普拉斯作用为
\[
  (\bm{L}\bm{x})[t]
  =
  2x[t]-x[(t-1)\bmod N]-x[(t+1)\bmod N].
\]
代入复指数模式$u_k[t]=N^{-1/2}e^{2\pi\mathrm{i}kt/N}$，
\begin{align*}
  (\bm{L}\bm{u}_k)[t]
  &=
  \left(
    2-e^{-2\pi\mathrm{i}k/N}
      -e^{2\pi\mathrm{i}k/N}
  \right)u_k[t]\\
  &=
  \left(2-2\cos\frac{2\pi k}{N}\right)u_k[t].
\end{align*}
因此DFT复指数也是环图拉普拉斯的特征向量。规则序列频率与图频率并非两个无关定义；
在环图上，它们由同一组模式连接。

这种对应也有边界。一般图的特征向量依赖节点连接和边权，不同图上的第$k$个谱系数
没有天然可比的物理频率。重特征值对应的特征子空间有多组正交基，单个特征向量还存在
符号或旋转不确定性。跨图比较时，应比较由子空间、谱区间或图构造共同确定的对象，
不能直接把两个图的同序号系数当作同一模式。

\subsection{谱图滤波与局部多项式是同一个算子}

给定标量函数$g$，可按矩阵函数定义图滤波器
\[
  g(\bm{L})
  =
  \bm{U}g(\bm{\Lambda})\bm{U}^{\top}.
\]
在图频域中，
\[
  \widehat{\bm{y}}_{\mathcal{G},k}
  =
  g(\lambda_k)\widehat{\bm{x}}_{\mathcal{G},k}.
\]
这与规则序列滤波器按频率响应逐点缩放完全对应。若$g$是$K$次多项式，
\begin{equation}
  \bm{y}
  =
  \sum_{r=0}^{K}\alpha_r\bm{L}^r\bm{x},
  \label{eq:signal-analysis-polynomial-graph-filter}
\end{equation}
则不必显式计算特征向量。一次乘$\bm{L}$只在一跳邻域内聚合差异，
$\bm{L}^r\bm{x}$只依赖至多$r$跳可达的节点，因此多项式谱滤波同时具有局部计算
解释。

例如三节点路径图的拉普拉斯与一个集中在首节点的信号为
\[
  \bm{L}=
  \begin{bmatrix}
    1&-1&0\\
    -1&2&-1\\
    0&-1&1
  \end{bmatrix},
  \qquad
  \bm{x}=
  \begin{bmatrix}1\\0\\0\end{bmatrix}.
\]
取一次多项式$g(\lambda)=1-\lambda/3$，则
\[
  g(\bm{L})\bm{x}
  =
  \left(\bm{I}-\frac{1}{3}\bm{L}\right)\bm{x}
  =
  \begin{bmatrix}2/3\\1/3\\0\end{bmatrix}.
\]
首节点的一部分数值只传到一跳邻居，没有越过第二个节点；在谱坐标中，同一计算把每个
系数乘以$g(\lambda_k)$。该图的特征值为$0,1,3$，相应响应为$1,2/3,0$，所以最高
图频率被消去。局部传播和谱缩放是同一个矩阵算子的两种读法。

这里的“卷积”是由选定图算子定义的谱滤波，不具备规则网格上的唯一核平移。
改变图构造、归一化拉普拉斯或边权会改变特征模式与滤波器；即使函数$g$的公式相同，
所得空间作用也可能不同。图信号方法保留了“在算子特征坐标中逐点缩放”的思想，
却不能无条件继承普通序列的平移与跨图频率含义。

\subsection{线性时不变状态递推产生卷积核}

卷积把每个输入位置对当前输出的贡献一次写完，状态递推则逐时保存过去信息。考虑
\begin{align}
  \bm{s}_{t+1}
  &=
  \bm{A}\bm{s}_t+\bm{B}\bm{x}_t,
  \notag\\
  \bm{y}_t
  &=
  \bm{C}\bm{s}_t+\bm{D}\bm{x}_t,
  \label{eq:signal-analysis-lti-state-space}
\end{align}
其中状态$\bm{s}_t\in\mathbb{R}^{d_s}$，输入
$\bm{x}_t\in\mathbb{R}^{d_x}$，输出$\bm{y}_t\in\mathbb{R}^{d_y}$，
矩阵$\bm{A}\in\mathbb{R}^{d_s\times d_s}$、
$\bm{B}\in\mathbb{R}^{d_s\times d_x}$、
$\bm{C}\in\mathbb{R}^{d_y\times d_s}$和
$\bm{D}\in\mathbb{R}^{d_y\times d_x}$不随$t$变化。递推一次、两次可以看到
\[
  \bm{s}_1=\bm{A}\bm{s}_0+\bm{B}\bm{x}_0,\qquad
  \bm{s}_2=\bm{A}^2\bm{s}_0+\bm{A}\bm{B}\bm{x}_0+\bm{B}\bm{x}_1.
\]
这里用$\bm{s}_t$表示状态、用$\bm{x}_t$表示输入，是为了延续本章把$x$作为信号的
约定。第~\ref{chap:dynamical-systems-and-state-estimation}章讨论同一个线性状态
过程时，改用$\bm{x}_k$表示状态、用$\bm{u}_k$表示输入；也就是说，
$\bm{s}_t$对应该章的$\bm{x}_k$，$\bm{x}_t$对应该章的$\bm{u}_k$，时间下标
$t$对应$k$。令该章的过程噪声为零，并取
$\bm{A}_k=\bm{A},\bm{B}_k=\bm{B}$，便得到这里相同的状态递推；这里的输出方程还
显式保留了直接传递项$\bm{D}\bm{x}_t$。
对一般$t$归纳得到
\begin{equation}
  \bm{s}_t
  =
  \bm{A}^t\bm{s}_0
  +
  \sum_{\tau=0}^{t-1}
  \bm{A}^{t-1-\tau}\bm{B}\bm{x}_{\tau}.
  \label{eq:signal-analysis-state-unrolling}
\end{equation}
归纳步骤只需把右侧乘$\bm{A}$并加$\bm{B}\bm{x}_t$，幂次与求和上界便同时向前
移动一格。

代入输出方程，
\begin{equation}
  \bm{y}_t
  =
  \bm{C}\bm{A}^t\bm{s}_0
  +
  \bm{D}\bm{x}_t
  +
  \sum_{\tau=0}^{t-1}
  \bm{C}\bm{A}^{t-1-\tau}\bm{B}\bm{x}_{\tau}.
  \label{eq:signal-analysis-state-convolution}
\end{equation}
若初始状态为零，定义矩阵值脉冲响应
\begin{equation}
  \bm{K}_0=\bm{D},\qquad
  \bm{K}_r=\bm{C}\bm{A}^{r-1}\bm{B}\quad(r\geq1),
  \label{eq:signal-analysis-state-kernel}
\end{equation}
便有
\[
  \bm{y}_t
  =
  \sum_{\tau=0}^{t}\bm{K}_{t-\tau}\bm{x}_{\tau}.
\]
这正是因果卷积：$\bm{K}_r$描述$r$步之前的输入经过状态传播后对当前输出的作用。
非零初态产生额外的自由响应$\bm{C}\bm{A}^t\bm{s}_0$，它不是输入卷积的一部分。

标量递推
\[
  s_{t+1}=as_t+bx_t,\qquad y_t=cs_t+dx_t
\]
给出$K_0=d$以及$K_r=ca^{r-1}b$。当$|a|<1$时，核随滞后几何衰减；
$|a|>1$时，远期输入的系数反而增长。这个观察只展示标量情形，矩阵系统的长期行为还
取决于特征值、非正规性和输入作用。相关稳定性条件将在
第~\ref{chap:dynamical-systems-and-state-estimation}章系统建立，本章不把一个形式展开
当作稳定保证。

取$a=1/2$、$b=c=1$、$d=0$、$s_0=0$，并把贯穿输入依次送入
$x_0=1,x_1=2,x_2=1$。逐步递推得到
\[
  y_0=0,\qquad y_1=1,\qquad y_2=\frac{5}{2},\qquad y_3=\frac{9}{4}.
\]
另一方面，状态核为$K_0=0,K_1=1,K_2=1/2,K_3=1/4$，所以
\[
  y_3=K_3x_0+K_2x_1+K_1x_2
  =\frac{1}{4}+1+1=\frac{9}{4}.
\]
递推与卷积给出相同结果，但输出从$y_0$开始包含一个由严格因果状态造成的零项；
若$d\neq0$，当前输入还会通过$K_0=d$立即进入输出。

\subsection{时变与输入依赖更新不再对应固定卷积}

若状态矩阵随时间改变，
\[
  \bm{s}_{t+1}
  =
  \bm{A}_t\bm{s}_t+\bm{B}_t\bm{x}_t,
\]
并令$\bm{y}_t=\bm{C}_t\bm{s}_t+\bm{D}_t\bm{x}_t$，
从时刻$\tau$传播到$t$的系数变成有序乘积
\[
  \bm{\Phi}(t,\tau+1)
  =
  \bm{A}_{t-1}\bm{A}_{t-2}\cdots\bm{A}_{\tau+1}.
\]
当$t=\tau+1$时，这个乘积按约定为单位矩阵。输出对输入的核依赖两个绝对时刻，
而不只依赖滞后$t-\tau$：
\[
  \bm{K}_{t,t}=\bm{D}_t,\qquad
  \bm{K}_{t,\tau}
  =
  \bm{C}_t\bm{\Phi}(t,\tau+1)\bm{B}_{\tau}
  \quad(0\leq\tau<t).
\]
这仍是线性输入输出关系，却不再平移不变，因而不是一个固定核的卷积。矩阵乘法次序
也不能交换；较晚时刻的状态转移必须写在左侧。

若$\bm{A}_t$或门控系数本身依赖$\bm{x}_t$或$\bm{s}_t$，展开后的“核”还会依赖
当前输入轨迹。此时系统通常是非线性的，叠加原理失效。可以在固定轨迹附近线性化，
得到局部的时变响应，但它不能作为对所有输入都成立的卷积核。区分固定线性时不变核、
两时刻的线性时变核和输入依赖的动态作用，是从本章进入动力系统分析时必须保留的边界。

\section{本章小结}
\label{sec:signal-analysis-summary}

有序输入上的线性作用之所以可以得到特殊算法，是因为平移结构让同一个脉冲响应在不同
位置复用。线性与平移不变共同推出卷积；零延拓产生线性卷积，周期延拓产生循环卷积，
相关运算则改变核方向与复共轭约定。边界、翻转和裁剪属于算子定义，不能留给实现任意
决定。

同一次卷积可以用矩阵、多项式或频率坐标表示。多项式乘积把卷积写成系数乘法，模
$z^N-1$解释循环回绕；DFT在单位根处求值，并用正交复指数恢复系数。卷积定理把循环
卷积化为频域逐点乘法，线性卷积只有在补零长度不小于完整输出长度时才能使用这一结果。
FFT利用单位根递归，Toom--Cook利用一般求值插值，Winograd类方法利用余数分解和
中国剩余定理。三者都重新组织计算，却具有不同的变换成本、数值误差和适用边界。

频率表示还说明信息在哪里丢失。滤波器的频率响应给出各模式的缩放；采样会周期复制
频谱，副本重叠后便发生混叠。抗混叠滤波可以在下采样前限制频带，却不能保留已被滤除
的细节。图拉普拉斯谱把“变化模式”推广到不规则关系，但模式依赖具体图，不能在不同
图之间按序号直接对应。

线性时不变状态递推展开后产生矩阵值卷积核，核的第$r$项是输入经过$r$步状态传播后
对输出的作用。非零初态、时变矩阵和输入依赖更新分别引入自由响应、两时刻核与非线性
动态，不能继续套用固定卷积。第~\ref{chap:dynamical-systems-and-state-estimation}章
将从同一个状态过程出发，研究轨迹、离散化、稳定性、随机扰动和不完全观测下的状态估计。
